\documentclass[12pt]{article}
\usepackage[utf8]{inputenc}
\usepackage[T1]{fontenc}
\usepackage{hyperref}
\title{Aweer Language Resources for K-12 Curriculum Development}
\author{Compiled by Claude Sonnet 5.5 (claude-sonnet-5-5) for Andr\'as Kornai}
\date{}
\begin{document}
\maketitle

\section{Introduction}
Aweer (ISO 639-3 bob; Glottocode awee1242), also known as Boni, is a Cushitic (Afro-Asiatic) language of coastal eastern Kenya, spoken in the Lamu and Ijara areas, with a smaller community in southern Somalia. Figures vary: about 7,600 speakers in the 2009 census and about 8,000 per Ethnologue, both as cited on Wikipedia. It is written in the Latin alphabet where written at all. Many speakers have shifted to Orma, Somali or Swahili. ``Boni'' may be considered pejorative; ``Aweer'' is preferred. Search found very little material; this report lists only what was found.

\section{Primary Educational Resources}
No primers, textbooks or school programmes for Aweer were found in search results.

\section{Linguistic and Reference Materials}
\begin{itemize}
\item Heine, \emph{Boni Dialects} (1982), a dialect atlas and grammar sketch, Glottolog reference record: \url{https://glottolog.org/resource/reference/id/19143}
\item Wikipedia, Aweer language: \url{https://en.wikipedia.org/wiki/Aweer_language}
\item Endangered Languages Catalogue entry (no standard orthography listed): \url{https://www.endangeredlanguages.com/lang/3948}
\end{itemize}

\section{Religious and Community Resources}
\begin{itemize}
\item Joshua Project, Aweer in Kenya (people-group profile; mentions a Bible society website as a resource but no confirmed translation): \url{https://joshuaproject.net/index.php/people_groups/10942/KE}
\item Wikipedia, Aweer people: \url{https://en.wikipedia.org/wiki/Aweer_people}
\end{itemize}

\section{Digital and Audio Resources}
No audio, video, corpus or software resources were found in search results.

\section{Curriculum Development Context}
No published dictionary, standard orthography or Bible translation was found. A curriculum would first need community-led orthography and lexicon work, possibly in cooperation with Kenyan language bodies and SIL. Existing documentation consists mainly of Heine's 1982 work.
\end{document}